\documentclass[11pt]{article}
\usepackage[a4paper,margin=1in]{geometry}
\usepackage{fontspec}
\usepackage[boldfont=false]{xeCJK}
\setCJKmainfont{FandolSong-Regular.otf}
\usepackage{amsmath}
\usepackage{amssymb}
\usepackage{array}
\usepackage{booktabs}
\usepackage{graphicx}
\usepackage{adjustbox}
\usepackage{enumitem}
\setlist[itemize]{topsep=2pt,partopsep=0pt,parsep=0pt,itemsep=2pt,leftmargin=1.4em}
\setlist[enumerate]{topsep=2pt,partopsep=0pt,parsep=0pt,itemsep=2pt,leftmargin=1.6em}
\usepackage{parskip}
\usepackage{fvextra}
\fvset{breaklines=true,breakanywhere=true,fontsize=\small}
\setlength{\parindent}{0pt}

\begin{document}

\begin{center}
  {\LARGE\bfseries Electrochemistry}\\[0.35em]
  {\large Igcse Chemistry}
\end{center}
\vspace{0.6em}

\section*{Electrolysis}

\textbf{Electrolysis} 电解 is the breaking down (\textbf{decomposition} 分解) of an \textbf{ionic compound} 离子化合物 — when it is \textbf{molten} 熔融 or in \textbf{aqueous} 水溶液 solution — by passing an \textbf{electric current} 电流 through it.

It only works when the substance is molten or dissolved, because then the \textbf{ions} 离子 are free to move and carry the charge. A solid ionic compound cannot be electrolysed because its ions are locked in place.

\subsection*{The electrolytic cell}

The set-up is called an \textbf{electrolytic cell} 电解池. Two \textbf{electrodes} 电极 (solid conductors) dip into the \textbf{electrolyte} 电解质 — the molten or aqueous substance being broken down.

\begin{itemize}
\item The \textbf{anode} 阳极 is the positive ($+$) electrode.

\item The \textbf{cathode} 阴极 is the negative ($-$) electrode.

\end{itemize}

The electrodes are often \textbf{inert} 惰性 (they do not react), such as \textbf{platinum} 铂 or carbon/\textbf{graphite} 石墨.

\subsection*{What is formed at each electrode}

There is a simple rule:

\begin{itemize}
\item \textbf{Metals} 金属 or \textbf{hydrogen} 氢气 are formed at the cathode.

\item \textbf{Non-metals} 非金属 (other than hydrogen) are formed at the anode.

\end{itemize}

For aqueous solutions, the product can depend on whether the solution is \textbf{dilute} 稀 or \textbf{concentrated} 浓. Here are the three Core examples:

\begin{center}
\adjustbox{max width=\linewidth}{%
\begin{tabular}{lll}
\toprule
\textbf{Electrolyte} & \textbf{At the cathode ($-$)} & \textbf{At the anode ($+$)} \\
\midrule
molten lead(II) bromide, $\text{PbBr}_2$ & \textbf{lead} 铅 (silvery liquid) & \textbf{bromine} 溴 (red-brown vapour) \\
concentrated aqueous sodium chloride & hydrogen (bubbles of gas) & \textbf{chlorine} 氯气 (pale green gas) \\
dilute \textbf{sulfuric acid} 硫酸 & hydrogen (bubbles of gas) & \textbf{oxygen} 氧气 (bubbles of gas) \\
\bottomrule
\end{tabular}}
\end{center}

\subsection*{How the charge moves}

During electrolysis:

\begin{itemize}
\item \textbf{Electrons} 电子 move through the wires (the \textbf{external circuit} 外电路) from the power supply.

\item At the cathode, positive ions gain electrons. Gaining electrons is \textbf{reduction} 还原.

\item At the anode, negative ions lose electrons. Losing electrons is \textbf{oxidation} 氧化.

\item Inside the electrolyte, the ions move: positive ions go to the cathode and negative ions go to the anode.

\end{itemize}

You can write a \textbf{half-equation} 半反应式 for each electrode. For molten lead(II) bromide:

\[
\text{Pb}^{2+} + 2e^{-} \rightarrow \text{Pb} \quad (\text{cathode, reduction})
\]

\[
2\text{Br}^{-} \rightarrow \text{Br}_2 + 2e^{-} \quad (\text{anode, oxidation})
\]

\subsection*{Electrolysis of copper(II) sulfate}

The product at the anode depends on the electrode:

\begin{itemize}
\item With \textbf{inert} carbon electrodes: \textbf{copper} 铜 forms at the cathode and oxygen forms at the anode. The blue colour of the solution slowly fades as copper is removed.

\item With copper electrodes: copper forms at the cathode, while the anode itself dissolves into the solution. The blue colour stays the same. This is used to purify copper.

\end{itemize}

\subsection*{Electroplating}

\textbf{Electroplating} 电镀 means covering a metal object with a thin layer of another metal. This improves its appearance and its resistance to \textbf{corrosion} 腐蚀 (rusting and wearing away).

To electroplate an object:

\begin{itemize}
\item the object to be coated is made the \textbf{cathode};

\item the plating metal is made the \textbf{anode};

\item the electrolyte is a solution containing ions of the plating metal.

\end{itemize}

\section*{Hydrogen–oxygen fuel cells}

A \textbf{fuel cell} 燃料电池 uses hydrogen and oxygen to make electricity directly. The only chemical product is water.

\[
2\text{H}_2 + \text{O}_2 \rightarrow 2\text{H}_2\text{O}
\]

It is useful to compare a fuel cell with a normal petrol (gasoline) engine in a vehicle.

\begin{center}
\adjustbox{max width=\linewidth}{%
\begin{tabular}{lll}
\toprule
\textbf{} & \textbf{Hydrogen–oxygen fuel cell} & \textbf{Petrol engine} \\
\midrule
Main product & water only & carbon dioxide and \textbf{pollutants} 污染物 \\
Effect on air & clean & adds to air pollution \\
Fuel storage & hydrogen is hard and dangerous to store (flammable, needs high pressure) & petrol is easy to store \\
Source & hydrogen may be made using fossil fuels & made from crude oil \\
\bottomrule
\end{tabular}}
\end{center}

\textbf{Advantages} of the fuel cell: the only product is water, so it does not pollute the air, and it changes chemical energy into electricity efficiently. \textbf{Disadvantages}: hydrogen is hard to store and transport safely, and producing the hydrogen can still use energy from fossil fuels.


\end{document}
